\chapter{Introduction}
~\textcolor{red}{\cite{TODO}}
In recent times, we are seeing that skills and job requirements are evolving rapidly.
As the culmination of various effects, including the development of new technologies and evolving consumer preferences, the labour force must adapt in order to satisfy the new requirements.
Given the historic pace of labour demand shifts, traditional methods of forecasting these changes are not well-equipped to handle this blistering, and accelerating, rate of change.

By these so-called "Traditional" methods, I mean that which accompany government or firm-level analysis of the skills market. These depend upon data which are limited in the following ways:

\begin{enumerate}
  \item They use structured data
  \item They collect the data periodically (e.g. Annually / Quarterly)
\end{enumerate}

An important feature of these methods is their approach to modelling the dynamics influencing labour demand.
Of key importance is the impact of innovation, and how new technology is estimated to affect skill requirements.

This paper aims to investigate whether Open Source Software (OSS) innovation Signals, derived from unstructured text data, can meaningfully contribute to the forecast of AI-related skill requirements present in \jobposts{}.

To break this down, there are a few key ideas being explored:
\begin{enumerate}
  \item Can open source software data represent information about Innovation in the AI sector?
  \item Can AI-related skill requirements be modelled the textual data within \jobposts{}?
  \item Can we integrate the two data channels to produce a quantified measure of their relationships?
\end{enumerate}

Put simply, if we can measure how similar an OSS \repo{} (innovation-source) is to \jobposts{}, in terms of the skill-profile and concepts involved, then we can represent how this innovation-source has diffused over time (i.e. by the date of \jobpost{}).
By relating this to the temporal development of the innovation-source, as represented by some other measure of the OSS \repo{}'s history, we can the development-diffusion timeline of that particular innovation.

If true, this insight has tangible benefits for the practice forecasting skill demand in the AI sector, within which open-source software is a prolific source of innovation.
Moreover, \jobposts{} \& OSS data are high-frequency and accessible, meaning that more granular and fine-tuned predictions can be produced in effective real-time.

% \section{Thesis Argument}
% As a student of philosophy, I find that a strong start to any paper is to frame the content and argument clearly.

% The function of this is to provide a clear view of what is either assumed to be true, what is \emph{proven} to be true (as follows from the relevant literature), or what I expect to \emph{demonstrate}.
% Thereafter, the conclusion necessarily follows, and it can be seen that the function of this paper is to prove the hypotheses as so labelled.
% \input{figures/argument.tex}

% \begin{argument}
\premises{
  % -----------------------------
  % NLP and skills extraction
  % -----------------------------
  \item There exist validated NLP methods for extracting technical skills from job postings with high precision and recall \cite{zhang-etal-2022-skillspan}.
  \item Contextual embeddings (e.g., Sentence-BERT) map job postings and repository descriptions into a shared feature space where cosine similarity reflects semantic similarity, as validated on semantic textual similarity benchmarks \cite{reimers_2019_SentenceBERTSentence}.
  \item Technical skills and repository descriptions are expressed in relatively standardized terminology, enabling reliable extraction and semantic comparison \cite{esco2023taxonomy}.

  % -----------------------------
  % GitHub activity as a signal
  % -----------------------------
  \item Stars indicate repository visibility and perceived popularity, while forks indicate community engagement and experimentation with the repository \cite{dabbish2012social}.
  \item Sustained repository activity (commits, issues, pull requests) signals ongoing experimentation and adoption of new technologies \cite{borges_2016_UnderstandingFactors}.
  \item Early engagement with repositories may reflect emerging technology trends that are later reflected in labor market demand.

  % -----------------------------
  % Job postings as demand proxies
  % -----------------------------
  \item Job postings represent explicit employer demand for skills and thus serve as a proxy for labor market demand trends.
  \item Temporal trends in job postings correlate with broader industry hiring trends \cite{radlinski2011learning,zhang-etal-2022-skillspan}.
  \item Skill mentions in job postings are described with relatively consistent terminology, enabling robust extraction using NLP techniques \cite{esco2023taxonomy}.

  % -----------------------------
  % Embeddings & similarity
  % -----------------------------
  \item Cosine similarity between embeddings in the same feature space provides a quantitative measure of semantic similarity suitable for ranking items \cite{mikolov2013distributed, reimers_2019_SentenceBERTSentence}.
  \item Cosine similarity preserves approximate local neighborhoods: if two items are both highly similar to a shared reference embedding, they are likely to be semantically related to each other \cite{ethayarajh2019contextual}.
  \item Cosine similarity induces a partial ordering of job postings with respect to each repository: if $\cos(J_A, R) > \cos(J_B, R)$, then $J_A$ is considered more semantically aligned with $R$ than $J_B$, enabling temporal aggregation and trend detection.
  \item Aggregated similarity between repositories and job postings can be used to detect alignment between emerging technologies and labor market demand.
  \item Increases in repository activity may precede increases in job posting similarity, justifying temporal alignment and forecasting analysis.
}

\end{argument}

\section{Aims \& Objectives}

\subsection{Primary Aim}
To develop and evaluate an AI-related skill demand forecasting methodology, integrating innovation signals (from Open-Source Software in particular) and using job ads as a proxy for skill demand.

\subsection{Objectives}
\begin{enumerate}
  \item \textbf{Data Collection:} Collect \& collate data for both channels in:
        \begin{enumerate}
          \item Open-source Software (\gh{}): \emph{textual content \& historic measure of development}
          \item \jobposts{} (e.g. LinkedIn, Indeed) - \emph{content \& date-posted}
        \end{enumerate}
  \item \textbf{Representation:} Develop a high-quality representation method for integrating textual data from both channels, involving:
        \begin{enumerate}
          \item \textbf{Build a Baseline:} A weak, keyword-based mechanism.
          \item \textbf{Advanced Representation:} An intelligent and robust method for representing text as embeddings.
          \item \textbf{Experimentation:} Fine tune and alter the components of embedding model, evaluating against the baseline.
        \end{enumerate}
  \item \textbf{Time series:} Investigate Relationship of derived time series, looking at Granger causality values and Pearson correlations.
  \item \textbf{Discuss:} Critically evaluate the derived results, With reflections on the benefits and limitations of the approach.
\end{enumerate}

\section{Roadmap}

In the following chapter, I discuss related pieces of work and how their findings informed the paradigm for this project's implementation.
Given the domain intersectional nature of this paper, works that are "related" include those on the topics of economic theory, natural language processing and AI itself.

Following this, the implementation of this project is detailed. Firstly, this concerns how data was acquired and preprocessed.
Thereafter, the primary component of the pipeline is discussed - namely, the representation pipeline.
Within this, I demonstrate both how the baseline model and the representation model are built, discussing the involved components and how experimentation can be conducted to find optimal constructions.
The subsequent section demonstrates the downstream task of these representation results - in the context of time series analysis, the findings of our representation pipeline are presented alongside the data of \gh{} Stars.
Having demonstrated the function of the entire pipeline, this next section therefore details the results and rationale of several sets of experiments.

What thereafter follows is a critical evaluation of the process, discussing the project's limitations and its successes, as well as the opportunities for future work.
Finally, We discuss in the conclusion the outcomes and results of this paper, and their implications.
